\documentclass{article}

\usepackage{ctex}
\usepackage{amsmath}
\usepackage{amsfonts}
\usepackage{graphicx}
\usepackage{setspace}
\usepackage{amssymb}
\usepackage{latexsym}
\usepackage{amsthm, euscript, makeidx, color, mathrsfs}
\usepackage[numbers,sort&compress]{natbib}
\usepackage{hyperref}
\hypersetup{
    colorlinks=true,
    linkcolor=blue,
    urlcolor=cyan,
    citecolor=cyan
}

\usepackage{geometry}
\geometry{
	a4paper,
	left=2.5cm,
	right=2.5cm,
	top = 3cm,
	bottom = 2cm
}

\newtheorem {theorem}{Theorem}[section]
\newtheorem {lemma}[theorem]{{\bf Lemma}}
\newtheorem {corollary}[theorem]{{\bf Corollary}}
\newtheorem {proposition}[theorem]{{\bf Proposition}}
\theoremstyle{remark}
\newtheorem {remark}{{\bf Remark}}[section]
\newtheorem {conclusion}[remark]{Conclusion}
\newtheorem {condition}{{\bf Condition}}[section]
\theoremstyle{definition}
\newtheorem {definition}{{\bf Definition}}[section]
\newtheorem {example}{{\bf Example}}[section]
\theoremstyle{plain}
\numberwithin {equation}{section}
\newtheorem{assumption}{Assumption}
\numberwithin{assumption}{section}

\allowdisplaybreaks

\begin{document}

\title{神经网络与深度学习课程大纲}
\author{数据科学与大数据技术专业大三课程}
\date{ }

\maketitle


\section*{课程名称}

神经网络与深度学习

\section*{适用对象}

数据科学与大数据技术专业大三学生

\section*{教材}

邱锡鹏《神经网络与深度学习》

\section*{课程目标}

\begin{enumerate}
    \item 理解神经网络和深度学习的基本概念、原理和应用。
    \item 掌握常用的神经网络模型（如前馈神经网络、卷积神经网络、循环神经网络）及其优化方法。
    \item 能够应用深度学习技术解决实际问题，如图像分类、自然语言处理等。
    \item 了解深度学习的最新研究进展和前沿技术。
\end{enumerate}

\section*{课程大纲}

\subsection*{第一部分：机器学习基础}

\begin{enumerate}
    \item 机器学习概述
    \begin{itemize}
        \item 概率论基础回顾
        \item 机器学习核心定义与分类
        \item 机器学习四要素（数据、模型、学习准则、优化算法）
        \item 泛化与正则化
    \end{itemize}
\end{enumerate}

\subsection*{第二部分：基础模型}

\begin{enumerate}
    \item 线性模型
    \begin{itemize}
        \item 线性回归模型
        \item 分类问题、线性模型与损失函数
        \item Logistic回归与Softmax回归
        \item 感知器与支持向量机（SVM）
    \end{itemize}
    \item 前馈神经网络
    \begin{itemize}
        \item 神经网络概览与前馈计算结构
        \item 激活函数
        \item 反向传播算法
        \item 神经网络的非凸优化与梯度下降问题
    \end{itemize}
    \item 卷积神经网络
    \begin{itemize}
        \item 卷积的定义与操作机理
        \item 卷积神经网络结构组成与特征提取
        \item 典型模型与应用（含图像与文本数据应用）
    \end{itemize}
    \item 循环神经网络
    \begin{itemize}
        \item 循环神经网络概念与原理 
        \item 序列应用模式 
        \item 参数学习（BPTT算法）与长程依赖问题（梯度爆炸与消失）
    \end{itemize}
\end{enumerate}

\subsection*{第三部分：进阶模型}

\begin{enumerate}
    \item 网络优化与正则化
    \begin{itemize}
        \item 网络优化的特点分析
        \item 优化算法的改进与网络参数初始化
        \item 各类归一化方法与超参数优化策略
        \item 泛化与正则化方法（Dropout、权重衰减等）
    \end{itemize}
    \item 注意力机制
    \begin{itemize}
        \item 注意力分配机制与原理
        \item 人工神经网络的软性注意力、键值对注意力机制
        \item 自注意力模型及其应用
    \end{itemize}
    \item 外部记忆
    \begin{itemize}
        \item 记忆网络与结构化的外部记忆计算（神经图灵机等）
        \item 基于神经动力学的联想记忆（如Hopfield网络结构及其能量函数）
    \end{itemize}
    \item 无监督学习
    \begin{itemize}
        \item 聚类分析机理定性分析
        \item 基于样本距离的代表性方法（K均值法、层次聚类）
        \item 无监督特征学习
    \end{itemize}
	\item 生成模型 
	\begin{itemize}
		\item 生成模型概述与分类
		\item 生成对抗网络（GAN）的原理与训练方法
		\item 变分自编码器（VAE）的原理与应用
		\item 扩散模型的原理与应用
	\end{itemize}
\end{enumerate}

\subsection*{第四部分：大模型与前沿技术初探}

\begin{enumerate}
    \item Transformer 基础与大语言模型
    \begin{itemize}
        \item Transformer算法架构（Encoder与Decoder）
        \item 现代预训练语言模型（BERT、GPT系列）
        \item 生成式模型的结构细节分析与解码
        \item 本地部署开源模型
    \end{itemize}
    \item 神经网络求解微分方程
    \begin{itemize}
        \item 物理信息神经网络（PINNs）基本原理
        \item 神经网络求解常微分与偏微分方程
        \item 算子学习网络（DeepONet/Neural Operator）简介
    \end{itemize}
    \item 强化学习初步
    \begin{itemize}
        \item 强化学习核心概念与马尔可夫决策过程（MDP）
        \item 基于价值的方法（Q-Learning等）
        \item 深度强化学习（DQN与策略梯度方法）
    \end{itemize}
        \item 现代AI工具与前沿应用
        \begin{itemize}
            \item 智能辅助编程工具（如 Trae）
            \item 智能体（Agent）的概念与开发框架
            \item 前沿对话模型及其生态（如 DeepSeek 等）
        \end{itemize}
	\end{enumerate}

\section*{执行大纲}

以下按周次排列：

第9周劳动节，第16周端午节。

\begin{enumerate}
    \item 机器学习概述
    \item 概率论基础回顾 
    \item 线性回归模型
    \item 线性分类模型
    \item 前馈神经网络
    \item 实验课：线性模型的实现与应用
    \item 卷积神经网络
    \item 循环神经网络
    \item 实验课：前馈神经网络的实现与应用
    \item 实验课：卷积神经网络的实现与应用
    \item 实验课：循环神经网络的实现与应用
    \item 网络优化与正则化
    \item 注意力机制
    \item 大语言模型
    \item 实验课：网络优化与正则化方法的实现与应用
    \item 实验课：注意力机制的实现与应用
    \item 复习课
\end{enumerate}
	
\section*{教学方法}

\begin{enumerate}
    \item 理论讲解：通过课堂讲授，深入讲解神经网络和深度学习的核心概念、模型和算法。
    \item 实践操作：结合编程练习和实验，使用Python和深度学习框架（PyTorch）实现经典模型。
    \item 项目驱动：通过小组项目，学生可以选择一个实际应用问题进行深度学习模型的开发与优化。
    \item 论文阅读：引导学生阅读深度学习领域的经典论文和最新研究进展，培养科研能力。
\end{enumerate}

\section*{考核方式}

\begin{enumerate}
    \item 平时作业：包括编程练习和理论题目，占总成绩的30\%。
    \item 项目报告：考查学生建模并解决问题的能力，占总成绩的20\%。
    \item 期末考试：考察基础知识、基础模型、理论与应用的知识点50\%。
\end{enumerate}

\section*{参考资源}

\begin{enumerate}
    \item 教材：邱锡鹏《神经网络与深度学习》
    \item 在线课程：
    \begin{itemize}
        \item 邱锡鹏主讲线上课程：
        \item 斯坦福大学CS231n（计算机视觉中的深度学习）
        \item 斯坦福大学CS224n（自然语言处理中的深度学习）
    \end{itemize}
    \item 深度学习框架：
    \begin{itemize}
        \item 神经网络与深度学习课程资源: https://nndl.github.io/
        \item 《动手学深度学习》（李沐）：https://zh-v2.d2l.ai/
        \item PyTorch: http://pytorch.org
    \end{itemize}
    \item 学术会议：
    \begin{itemize}
        \item ICLR（国际表示学习会议）
        \item NeurIPS（神经信息处理系统年会）
        \item ICML（国际机器学习会议）
    \end{itemize}
\end{enumerate}

\end{document}